\section{Theorie}
\label{sec:Theorie}
Digitale Schaltungen verknüpfen Eingangs- und Ausgangssignale logisch über verschiedene Gatter. Dabei arbeiten sie mit den zwei Aussagen bzw. Zuständen „1“ (wahr) und „0“ (nicht wahr). Alle digitalen Schaltungen basieren auf Transistorschaltungen. Transistoren reagieren auf Spannungen und Ströme. Je nach verwendeter Halbleitertechnologie, zum Beispiel TTL oder CMOS, werden unterschiedliche Spannungswerte als High (H) oder Low (L) erkannt. Bei der positiven Logik steht H für 1 und L für 0. Bei der negativen Logik ist es genau umgekehrt, H steht für 0 und L für 1.\\
Operationen \\
Digitale Schaltungen lassen sich in statische und zeitabhängige Schaltungen unterteilen. Bei statischen Schaltungen ist der Zeitpunkt eines Zustands unerheblich. Bei zeitabhängigen Schaltungen hingegen verändert sich der Ausgang in Abhängigkeit von einer Änderung am Eingang. Dabei spielen sowohl der Zeitpunkt als auch die vorherigen Zustände (Vorgeschichte) eine wichtige Rolle.\\
Ein klassisches Beispiel für eine solche zeitabhängige Schaltung ist der RS-Flip-Flop. Er besteht im Kern aus vier NAND-Gattern und kann zwei stabile Ausgangszustände einnehmen (bistabile Kippstufe). Über die Eingänge S (Set = Setzen) und R (Reset= Rücksetzen bzw. Löschen) wird der Zustand des Flip-Flops gesteuert. An den Ausgängen gelten dabei komplementäre Zustände. Ist beispielsweise der Ausgang $Q=1$ gesetzt, so gilt für den Ausgang  $ \overline{Q} = 0$. Dieser Zustand bleibt so lange stabil gespeichert, bis ein Signal am R-Eingang eine Änderung auslöst.\\
RS-Flip-Flops werden häufig zur Speicherung von Daten verwendet. Dabei wird die Aufnahme und Ausgabe der Daten durch ein Taktsignal (Clock) gesteuert. Liegt am C-Eingang ein High-Pegel (H) an, übernimmt die Schaltung die am Dateneingang (D) anliegende Information. Liegt hingegen ein Low-Pegel (L) an, bleibt der bisherige Zustand gespeichert und unverändert.\\
Jedoch hat eine RS-Flip-Flop Schaltung auch seine Schwächen. Die Eingangskombination S=R=1 ist verboten, da die Kippstufe nicht gleichzeitig gesetzt und rückgesetzt werden kann. Dies führt zu einem undefinierten Zustand mit unvorhersehbarem Verhalten. Um den kritischen, verbotenen Zustand zu vermeiden, wird meist ein D-Flip-Flop verwendet. \\
Der D-Flip-Flop basiert auf einem RS-Flip-Flop, bei dem der Rücksetzeingang das entgegengesetzte Signal des Setzeingangs erhält \cite{DFlipFlop}. Er ist also ein taktgesteuerter RS-Flip-Flop in Verbindung mit einem NICHT-Gatter, das so geschaltet ist, dass immer S= $ \overline{R} $ gilt. Durch diese Konstruktion wird der R-Eingang nicht mehr von außen angesteuert, da das Eingangssignal direkt über das NICHT-Gatter invertiert weitergeleitet wird und somit entfällt. Die Ansteuerung selbst kann auf verschiedene Arten erfolgen. Beispielsweise ist eine statische Ansteuerung über einen Spannungsimpuls möglich. Eine dynamische Ansteuerung über eine steile Spannungsflanke, ist eine weitere Möglichkeit.\\
Legt man an den Takt-Eingang eine periodische Rechteckspannung an, halbiert sich die Ausgangsfrequenz, wodurch das Flip-Flop als Frequenzteiler wirkt. Bei n zusammengeschalteten Flip-Flops wird die Frequenz entsprechend durch $2^n$ geteilt.\\
Verbindet man Flip-Flops als Binärzähler, durchlaufen sie definierte Zahlenwerte. Ein 3-Bit-Zähler zählt beispielsweise von 0 bis 7, während vier Flip-Flops das Zählen bis 15 ermöglichen. Durch zusätzliche Gatter lässt sich ein Zähler beim Erreichen eines bestimmten Wertes auf Null zurücksetzen. Ein Zähler, der von 0 bis 9 zählt und dann neu startet, wird als BCD-Zähler bezeichnet.\\
Zudem lassen sich Signalfrequenzen messen, indem man die Impulse zählt, die innerhalb einer bestimmten Torzeit ankommen. Für die Erzeugung dieser Torzeit verwendet man häufig einen Quarzoszillator, dessen Frequenz über Teiler herabgeregelt wird.\\